\section{Conclusion: A Calculus and Its Vacuum Test}

The general rules of Nested Persistence now have independent mathematical formulations. A feasible division services its stated demand within total capacity including overhead; a reversible encoding retains the distinctions lost by its coarse description; an interface operates within an explicit reflected-load budget; and finite nesting depth follows from a positive per-level floor or a decreasing discrete resource. Lean checks these results without importing the $E_8$ constants example. Feasibility is distinguished from saturation, maximal gain and exact matching.

\subsection{What the Constants Example Establishes}
The vacuum application supplies proposed operational ledgers and response laws. Under those premises, Lean verifies the overhead and gauge-allocation arithmetic, scalar resource bounds, conditional charge selection, mass relations, and several certified numerical intervals. These implications are useful because their hypotheses and possible failures are explicit. They are not a proof that every persistent substrate must use the same ledgers or response laws.

The joint checks are decisive for the next stage, and several now point to a definite reading. The VEV fixes the electron Yukawa through the mass law; the earlier independent $\sigma/D$ dressing is excluded, and no normalization choice could have repaired it. The W/Z ratio holds at every interface scale, but the absolute masses need one electroweak interface factor that the strong readout cannot share; with the measured Z mass the proposal predicts $m_W\approx80.385$~GeV and $\Delta r\approx0.0351$, distinct from standard radiative corrections on the same inputs. At the scalar interface, asking what persists selects the largest gain that pays for its own reflection, which lies just below $128/989$. The strong-coupling running test must use the flavour thresholds implied by Paper~V's own masses. The core unity condition is exact by its Planck-mass definition. None of these distinctions is removed by a numerical goodness-of-fit summary.

The gauge calculations share a conditional event-law response and a common calibration. On a supplied required-resource set with positive additive refresh costs, the list model has one refresh per distinct resource at its minimum. Actual carrier operations now reproduce the eight- and twelve-field support columns, but those columns do not by themselves establish the physical resource assignments. The proposed ledgers remain $43$, $129+5+16\eta$ and $12\eta$. Under the intensive relational-density interpretation and a uniform carrier-based reference-direction sampler, the colour factor follows without inserting $\eta$ into the load definition. A fixed-direction reference history does not give that factor. The sampler, primitive-work interpretation and common readout/repair normalization therefore remain explicit obligations; numerical cost inequalities alone do not select the intensive contract.

The W/Z link's mathematical response has been strengthened independently of those physical choices. Its explicit generator is Hermitian, so Lean constructs a spectral unitary transport and a positive small-step radius, with a controlled fourth-order remainder giving the stated mass form (\texttt{LinkSpectralResponse}). No eigenbasis needs to be supplied as a premise. The current numerical entry point, \texttt{CurrentConstants}, separates the colour-ledger strong coupling and VEV-defined electron Yukawa from the earlier formulas, which remain as regression and exclusion results. Neither step supplies the absolute mass scale or the electroweak interface response.

\subsection{The Remaining Substrate Derivation}
The aim remains to derive physical constants from the least costly structure that can persist. To reach that conclusion, the workload and reference cost, spectral interpretation of the encoding floor, response normalization and scale dependence must follow from the substrate. In particular, the proposed beta coefficients, flavor invariant, strong-CP argument and particle-exclusion claims require their own dynamical or representation-theoretic proofs. The common representation adds a specific normalization test: its packet indices are $2,2,10/3$ for weak, colour and $Y/2$. Reading the proposed ledgers as raw refresh weights under one uniform packet readout gives $5547/37777$, rather than the proposed weak fraction. A shared calibration cannot remove that difference. Interpreting the ledgers as already normalized response weights therefore requires an explicit operational conversion; the derived compatibility factor is not an input to choose for agreement. The current formalization does not establish a solution to the hierarchy, strong-CP or vacuum-energy problems.

A proposed symmetric precursor illustrates the distinction between arithmetic and physical matching. The ratio $(\sigma-\chi)/(2D)=3/8$ follows from the given integers. Reading it as a unification-scale weak angle, or deriving the low-energy allocation $\Delta/(D\Delta+\nu\eta+\sigma)$ from it, requires an explicit response and scale-evolution map. That map has not been proved here.

\subsection{Next Steps}
The next tasks should follow their shared dependencies. The common actions and controlled comparison response are now constructed on the actual packet. The existing weak and colour gates have six carrier-type orbits, and the weighted response identifies the precise preparation constraints needed for the proposed ratios. Derive those preparation/work weights and one primitive-event workload contract for all three ledgers; then close scalar regulation, the mass scale and the interface response. A candidate reference walk should be constructed and priced when that contract requires it, rather than adopted solely to recover a target average. Sector exchange, generation classification and strong CP require an actual update symmetry and their remaining spectral and phase conditions. Running, gravity and neutrino propagation then need the corresponding companion dynamics. Persistence Mechanics supplies constraints to be tested across these stages; the constants example remains a conditional quantitative program rather than a completed uniqueness theorem for nature.

\begin{acknowledgments}
The author is an independent researcher and received no external funding for this work.

I would like to thank Brian Sheppard for rigorous and constructive feedback.

I would like to thank my friends and family for their patience and support throughout decades of discussions as I tried to understand every field I became interested in and the overarching pattern I kept seeing.
\end{acknowledgments}
